\par\noindent\begin{minipage}{\linewidth}
\small
\setlength{\tabcolsep}{4pt}
\renewcommand{\arraystretch}{1.13}
\captionof{table}{Estabilidad de las propiedades geométricas ante la selección.}\label{tab:S12}
\begin{tabular}{@{}>{\raggedright\arraybackslash}p{3.5cm}>{\raggedright\arraybackslash}p{3.4cm}>{\raggedright\arraybackslash}p{1.2cm}>{\raggedright\arraybackslash}p{1.5cm}>{\raggedright\arraybackslash}p{2.5cm}>{\raggedright\arraybackslash}p{2.6cm}@{}}
\toprule
\textbf{Propiedad} & \textbf{Diseño de selección} & \textbf{Casos} & \textbf{Alertas} & \textbf{Amplitud mediana} & \textbf{Amplitud máxima} \\
\midrule
Apertura angular & 20 medias muestras & 260 & 0 & 0,79\% & 1,81\% \\
Apertura angular & 5 selecciones de igual tamaño & 260 & 0 & 0,49\% & 1,85\% \\
Dimensión lineal efectiva & 20 medias muestras & 260 & 4 & 4,06\% & 12,84\% \\
Dimensión lineal efectiva & 5 selecciones de igual tamaño & 260 & 0 & 2,60\% & 7,53\% \\
Conexión & 20 medias muestras & 260 & 258 & 18,62\% & 423,07\% \\
Conexión & 5 selecciones de igual tamaño & 260 & 51 & 10,59\% & 289,69\% \\
\bottomrule
\end{tabular}
\end{minipage}\par
\par\smallskip{\small Un caso combina modelo y área. La amplitud es el rango central del 90\% dividido por el valor principal; también se comprueba el desplazamiento de la mediana. Umbrales: 5\% apertura, 10\% PR, 20\% conexión. Son umbrales operativos de sensibilidad, no normas de la revista ni precisión poblacional. Las cuatro alertas de PR en medias muestras siguen visibles aunque pasen las cinco selecciones de igual tamaño. PR es dimensión lineal efectiva, no cantidad de temas. Las alternativas espectrales comparten el espectro de covarianza y no validan independientemente el significado temático.\par}
